\section{The computation}\label{sec:method}

\subsection{Commutation as a finite constraint}
Let $g$ be a block map of radius $R$ with local rule $\gamma$, and $F=F_f$. Both $g\circ F$ and $F\circ g$ are block maps of radius $R+1$, so they are equal if and only if their local rules agree on every window $w\in\A^{2R+3}$. Write $n=2R+1$ and, for a window $w=(w_0,\dots,w_{2R+2})$, let $e_1,e_2,e_3$ be the three subwords of length $n$ starting at positions $0,1,2$ and let $e_L=F(w)$ be the word of length $n$ obtained by applying the local rule of $F$ to $w$. Then
\[
(g\circ F)(w)=\gamma(e_L),\qquad (F\circ g)(w)=\Bigl(\gamma(e_2)_y,\ \gamma(e_2)_x\oplus f\bigl(\gamma(e_1)_y,\gamma(e_2)_y,\gamma(e_3)_y\bigr)\Bigr),
\]
where $\gamma(e)_x,\gamma(e)_y$ are the two output bits. Writing $X_e,Y_e$ for the two unknown bits of table entry $e$, each window contributes the two constraints
\begin{equation}\label{eq:constraints}
X_{e_L}=Y_{e_2},\qquad Y_{e_L}\oplus X_{e_2}=f(Y_{e_1},Y_{e_2},Y_{e_3}).
\end{equation}
The first is two binary clauses; the second is a constraint on at most five variables, encoded by its truth table in at most sixteen clauses (the indices $e_L,e_1,e_2,e_3$ may coincide, and the encoding is over the set of distinct variables). Altogether there are $2\cdot4^{2R+1}$ variables and at most $18\cdot4^{2R+3}$ clauses; for rule $110$ there are $17{,}680$ clauses at $R=1$, $4.7\times10^6$ at $R=3$ and $7.5\times10^7$ at $R=4$.

\subsection{Enumeration and verification}
All solutions are enumerated by repeated solving with blocking clauses, using PySAT \cite{PySAT} with CaDiCaL~1.5.3 \cite{CaDiCaL} for $R\le3$ and MiniSat~2.2 \cite{MiniSat} at $R=4$, where CaDiCaL's memory use exceeded the machine. Every table returned by the solver is re-checked, independently of the CNF, by evaluating $g\circ F$ and $F\circ g$ on all $4^{2R+3}$ windows in NumPy. Each solution is then classified: it is matched against the $\sh^aF^b$ with $|a|+|b|\le R$, which are built directly by composing local rules; the two constants are recognised; everything else is labelled \emph{other}. The program asserts that every $\sh^aF^b$ with $|a|+|b|\le R$ occurs among the solutions, which checks the encoding against the only maps known in advance to commute.

Every solution is thus verified independently of the solver. Exhaustiveness is of another kind: the clauses are necessary conditions for commutation, so every commuting block map is a solution, and that there is no solution beyond those found rests on the final answer of the solver that the clauses, together with the blocking clauses of the solutions found, are unsatisfiable. This answer was cross-checked in two ways (\texttt{code/census/crosscheck.py}, Appendix~\ref{app:census}). First, every enumeration was repeated with other solvers: for $R\le3$ with Glucose~4 \cite{Glucose}, MiniSat~2.2 and once more CaDiCaL~1.5.3, on all $200$ instances that finish below the cap; for $R=4$ with Glucose~4 and CaDiCaL~1.5.3 for the rules $110$ and $30$, and with Glucose~4 for the other $39$ rules that finish below the cap. In every instance the solvers return the same set of tables, and this set agrees with the log of the enumeration. Second, for $R\le3$ the unsatisfiable formula was given to Glucose~4 with proof logging, and the clausal proof it emits was checked by a short forward reverse-unit-propagation checker that uses no SAT solver; this was done for $197$ of the $200$ instances, among them all those behind Theorem~\ref{cl:thm:census} and the cases $R\le3$ of Theorem~\ref{thm:110}, the exceptions being three instances with more than $1700$ solutions. No proof was checked at $R=4$, where exhaustiveness rests on the agreement of two or three solvers.

\subsection{Invertibility}
A block map is an automorphism of $\A^{\Z}$ if and only if it is injective. Injectivity of a block map of radius $R$ is decided exactly by the pair graph: its vertices are the pairs $(u,v)$ of words of length $2R$, with an edge $(u,v)\to(u',v')$ labelled $(a,b)\in\A^2$ whenever $u'=ua$ and $v'=vb$ after dropping the first letter, and $\gamma(ua)=\gamma(vb)$. Then $g$ fails to be injective if and only if a bi-infinite path passes through an off-diagonal vertex, which is decided by iteratively deleting vertices without a predecessor or without a successor. We implemented this twice, in pure Python and vectorised, and asserted agreement at $R=1$; at $R\le2$ every map found injective was also checked to be injective on all spatially periodic configurations of period at most $2R+4$.

For $R=2$ with rule $54$, where the full centralizer is too large to enumerate, we enumerated only the invertible part, as follows. An invertible block map is surjective, and a surjective block map of radius $R$ is \emph{balanced} \cite[Theorem~5.4]{Hedlund1969}: each output symbol occurs exactly $4^{2R}$ times in the table. Balance is added to the CNF as four cardinality constraints, the balanced commuting maps are enumerated, and each is submitted to the pair-graph test. Balance is necessary and the pair-graph test is exact, so this yields $\Cent^{\times}_R(F)$ exactly.

\subsection{Linear rules}
When $f$ is linear over $\F_2$, the constraints \eqref{eq:constraints} are homogeneous linear equations in the table bits and $\Cent_R(F)$ is an $\F_2$-subspace of table space. Its dimension was computed by Gaussian elimination: the equalities $X_{e_L}=Y_{e_2}$ are merged by union--find, which reduces the $2\cdot4^{2R+1}$ bits to $4^{2R}$ classes, and the remaining equations are reduced by bit-packed elimination. For rule $90$ and $R\le2$ the result was compared with the SAT enumeration of all $2^{\dim}$ solutions; this tests the encoding and the solver on a control whose solution set is known independently.

